% ============================================================================
% CHAPITRE 3 : CONCLUSION
% ============================================================================

\section{CONCLUSION DU CHAPITRE~\ref{chap:constructions}}
\label{sec:conclusion_constructions}

\begin{sloppypar}
La famille SOTP (section~\ref{sec:sotp_construction}) propose trois
constructions de r\'ef\'erence, SOTP-Low, SOTP-Base et SOTP-High, adapt\'ees
chacune \`a un r\'egime de richesse contextuelle. Leurs propri\'et\'es
cryptographiques, \`a savoir la confidentialit\'e (calculatoire pour les
variantes \`a PRG, divulgation nulle de connaissance pour SOTP-High) et la
fiabilit\'e de v\'erification, sont prouv\'ees sous PRG, DLog et q-SDH/q-PKE. Les facteurs
$0{,}75$, $0{,}5$ et $0{,}45$ des formulations initiales sont conserv\'es comme
heuristiques de conception, tandis que l'optimalit\'e universelle
est formul\'ee comme conjecture (Conjecture~\ref{conj:optimalite_universelle}). Cette
r\'eserve a une raison structurelle, car une couche Clay qui ne d\'epend que du contexte public
n'apporte aucune information conditionnelle au sens de $\Opp$
(Remarque~\ref{rem:statut_sotp}).

La v\'erification Lean4 (section~\ref{sec:lean4}) a certifi\'e
le noyau alg\'ebrique. Les in\'egalit\'es de Pinsker et de Fano y sont
int\'egr\'ees comme propri\'et\'es des instances de protocole et les hypoth\`eses
institutionnelles y sont explicit\'ees~; la borne
$\Delta/(\gamma\lambda) + O(\lambda^{-2})$ et sa forme raffin\'ee sont v\'erifi\'ees
sans \texttt{sorry}. Deux t\'emoins formalisent la coh\'erence (une instance du
mod\`ele abstrait atteint la borne rel\^ach\'ee) et la port\'ee (une instance hors
r\'egime atteint $1/4$). Cette d\'emarche illustre la valeur
\'epist\'emique de la v\'erification formelle, qui certifie des \'enonc\'es et contraint aussi
\`a expliciter ce qui \'etait rest\'e implicite. Elle a conduit \`a des formulations
plus pr\'ecises de trois \'enonc\'es publi\'es (sous-section~\ref{subsec:lean4_erreurs}).

Quatre limitations demeurent. Les param\`etres
$\Delta$, $\alpha$, $\beta$ reposent sur des estimations analytiques. Leur mesure directe aupr\`es
d'acteurs institutionnels r\'eels rel\`eve des protocoles du
chapitre~\ref{chap:experimentation}. Les m\'etriques CLO
des syst\`emes de r\'ef\'erence ont \'et\'e fix\'ees par calibration param\'etrique
r\'etrospective, sans exp\'erimentation directe et ne valident pas les
heuristiques SOTP. La formalisation Lean4 couvre le noyau
alg\'ebrique. Les fondements probabilistes et les r\'eductions de s\'ecurit\'e
des constructions SOTP restent hors de son p\'erim\`etre. Le mod\`ele quantique utilise
principalement le canal d\'epolarisant, qui peut ne pas capturer toutes les
nuances des syst\`emes quantiques r\'eels.

La cryptographie institutionnelle dispose d\'esormais des m\'etriques, de la
loi d'impossibilit\'e, des constructions de r\'ef\'erence et de la certification
formelle que requiert tout sous-domaine th\'eorique. Ces r\'esultats doivent
encore \^etre confront\'es aux syst\`emes r\'eels, ce qui conduit \`a se demander
si les param\`etres institutionnels qui gouvernent ces bornes sont mesurables.
Cette question correspond \`a l'\'etape \textbf{X} de la progression et fait
l'objet du chapitre~\ref{chap:experimentation}.
\end{sloppypar}
